\chapter{Chemistry: Counting Reactions}
\label{chap:chemistry-counting-reactions}

Chemistry is where the instrument's count becomes visible to the eye. Everywhere else the count is inferred
--- read off a trajectory, a field, a spectrum --- but in chemistry it is the plain content of the science:
reactions balance in whole numbers, compounds form in fixed integer proportions, and the continuum of masses
and temperatures is only an alphabet for stating integer facts. This is counting at its purest, the
instrument working in the one place where the world hands it whole numbers directly, because its units ---
atoms --- are indivisible and cannot be had in fractions. Four effects carry the chapter: definite
proportions, the integer balance of a reaction, the arbitrary-but-stable temperature scale, and the catalyst
that changes a reaction's path without changing its count.

\section{The Archimedes--Proust Effect: definite proportions}
\label{se:archimedes-proust}

At the turn of the nineteenth century two chemists argued about whether a compound's composition was fixed or
free. Berthollet held that two elements could combine in any proportion across a continuous range; Proust held
that a given compound always contained its elements in the same fixed ratio --- water always eight parts
oxygen to one of hydrogen by mass, however it was made. Proust was right, and the reason is the instrument's:
there are two modes of measurement, the continuous, where geometry supplies arbitrarily fine values, and the
integer, where commitment fixes a whole number, and a compound belongs to the second. Atoms are indivisible
counting units, and a compound is a whole-number ratio of them. Water is two hydrogen atoms to one oxygen,
always, and no preparation, however careful, yields two-and-a-half hydrogens to an oxygen, because there is no
half an atom to count.

The honest floor is a finite count obligation: the proportions are integers, the integer balance is the fact,
and the continuous masses are alphabet, not content. The reading is that a compound is an integer relation.
The continuum can express it to any precision --- sixteen grams of oxygen to two of hydrogen, eight to one ---
but the fact it expresses is a count: two of the one, one of the other. This is chemistry's first count, and
it is the bare fact that atoms come in whole numbers, made into a law.

Dalton sharpened the point into the strongest early evidence for atoms. The same two elements can form more
than one compound --- carbon and oxygen make both carbon monoxide and carbon dioxide --- and when they do, the
masses of the one element combining with a fixed mass of the other stand in small whole-number ratios. Hold
the carbon fixed, and the oxygen in the dioxide is exactly twice the oxygen in the monoxide: one to two, not
one to one-point-nine. This is the law of multiple proportions, and it is nearly a proof on its own that
matter is made of indivisible units. A continuum would have no reason to prefer the ratio two; counting units
have no choice but to.

\section{The Stoichiometry Effect: a reaction is an integer program}
\label{se:stoichiometry}

If a compound is an integer ratio, a reaction is an integer balance. Burn hydrogen in oxygen and you get
water --- but not in any proportion: the equation must balance, atom for atom,
\begin{equation}
  2\mathrm{H}_2 + \mathrm{O}_2 \to 2\mathrm{H}_2\mathrm{O},
\end{equation}
two hydrogen molecules and one of oxygen making two waters, every atom on the left accounted for on the right.
In general a merge of $N_A$ and $N_B$ into $N_C$ is admissible only if there are integers $a, b, c \in
\mathbb{Z}$ that balance every atom exactly,
\begin{equation}
  a\,N_A + b\,N_B = c\,N_C,
\end{equation}
and no fractional event may be recorded --- no half a molecule reacts, no partial refinement is admitted. This
is Lavoisier's conservation made into arithmetic: nothing is created or destroyed in the reaction, so every
atom that enters must leave, and the balance is an integer bookkeeping the instrument can check by counting.

The balance has a sharp linear-algebraic form, and naming it ties the reaction to the relation-algebra the
book uses elsewhere. Build the formula matrix $A$: one row per element, one column per species, each entry the
number of atoms of that element in that species. A reaction is a vector $\mathbf{v}$ of stoichiometric
coefficients --- positive for products, negative for reactants --- and conservation, the demand that every
atom entering leaves, is the single statement that $A$ annihilates $\mathbf{v}$,
\begin{equation}
  A\mathbf{v} = \mathbf{0}, \qquad \mathbf{v} \in \mathbb{Z}^n .
\end{equation}
Balancing a reaction is finding a vector in the kernel of the formula matrix, and the demand that $\mathbf{v}$
be integer, not merely real, is the demand that atoms come in whole units. The kernel as a linear subspace
would admit any scalar multiple, a continuum of solutions; the integer kernel admits only whole-number ratios,
and the smallest such vector --- the coefficients in lowest terms --- is the balanced equation. An equation
that balances is one whose formula matrix has a nonzero integer null vector; an unbalanceable one is a matrix
whose kernel meets the integer lattice only at zero.

The honest floor is a finite count obligation: the reaction is exactly the integer balance, and admissibility
is solvability in whole numbers. The reading is that a chemical reaction is an integer program. It proceeds if
and only if the atoms balance in whole numbers, and the selection rule governing which reactions can happen at
all is just the demand that an integer solution exist --- an unbalanceable equation is not a slow reaction or
an inefficient one; it is simply not a reaction.

\section{The Celsius--Lagrange Effect: a scale arbitrary in symbol, stable in relation}
\label{se:celsius-lagrange}

Not everything in chemistry is a count; some of it is a scale, and a scale is a different kind of object,
built to support interpolation without claiming the underlying phenomenon is continuous. Celsius fixed two
reproducible anchors --- the freezing and boiling of water, called nought and a hundred --- and interpolated
linearly between them,
\begin{equation}
  T = T_0 + (T_1 - T_0)\,\frac{x - x_0}{x_1 - x_0},
\end{equation}
so that any intermediate temperature is recoverable by the rule. But the scale is arbitrary in its symbols and
stable only in its relations. The anchors are facts and the relations among readings are facts; the off-anchor
values are scaffolding --- a convention chosen to interpolate, not a claim that the phenomenon is genuinely
continuous there. That the numbers are free is plain from the rival scales: Fahrenheit and Kelvin set
different anchors and different zeros, and a temperature reads as a different number on each, yet all three
encode the same orderings and the same ratios of intervals. What every scale agrees on is the relation; what
each is free to choose is the label.

The honest floor is a smooth-shadow analogy: the continuous scale is the smooth shadow of two anchors and a
rule, and the bridge from the anchors to the continuum of readings is an interpolation, named as one. The
reading is that a scale's numbers are arbitrary but its relations are not. What a thermometer reports between
its anchors is scaffolding built to interpolate; the only facts are the anchors and the order and ratios the
scale preserves. The symbol is free; the relation is fixed.

\section{The First Effect of Gibbs: a catalyst changes the path, not the count}
\label{se:gibbs-catalyst}

A catalyst lowers the strain a reaction must overcome to proceed, without altering the net ledger. Introduce a
catalytic structure $K$ and it supplies an alternate sequence of intermediate steps that reduce the curvature
of the admissible path --- an easier route over the same terrain. But $K$ is regenerated: it appears in the
forward step and is returned by the inverse, so the forward and inverse catalyst contributions cancel and the
net balance is exactly what it was without $K$. The count stays additive. The Haber process is the standing
example: an iron catalyst lets nitrogen and hydrogen combine into ammonia at a workable rate, and the iron is
returned untouched at the end, nowhere in the net equation $\mathrm{N}_2 + 3\mathrm{H}_2 \to 2\mathrm{NH}_3$.
It changed how fast the reaction ran and by what path, but not a single integer in the balance.

The honest floor is a finite count obligation: the catalyst cancels in the net count, leaving the integer
balance of the reaction untouched. The reading is that a catalyst changes the path, not the destination. It
lowers the strain of getting there and is handed back unchanged, so the reaction's count --- its
stoichiometry, its definite proportions --- is exactly what it would be without it. The catalyst speeds the
reconciliation; it does not alter the books.

\bigskip

Chemistry, read off finite records, is integer arithmetic made visible: definite proportions, reactions that
balance in whole numbers, scales whose symbols are free but whose relations are fixed, and catalysts that
change the path without touching the count. Nowhere did the instrument need the continuum; the atoms handed it
whole numbers, and it counted them. Part III closes next with condensed matter, where these counts settle into
the states of solids, and the instrument meets the last of its four honest floors.
